\documentclass[10pt,a4paper]{amsart}
\usepackage[margin=19mm]{geometry}
\usepackage{amsmath,amssymb,mathtools}
\usepackage[scr=rsfs,cal=euler]{mathalfa}
\usepackage{xfrac}
\usepackage[unicode,hidelinks,bookmarks=false]{hyperref}
\newcommand{\ie}{i.e.\ }
\newcommand{\cf}{cf.\ }
\newcommand{\resp}{respectively\ }
\newcommand{\Subsection}{\subsection}
\providecommand{\nobreakdash}{\nobreak\hbox{-}}
\setlength{\emergencystretch}{2em}
\multlinegap0pt
\usepackage{amssymb}
\usepackage{xcolor}
\usepackage[shortlabels]{enumitem}
\usepackage{mathtools}
\newtheorem{lemma}{Lemma}
\newtheorem{proposition}{Proposition}
\newcommand{\R}{\mathbb{R}}\newcommand{\C}{\mathbb{C}}
\newcommand{\cP}{\mathcal{P}}
\newcommand{\cS}{\mathcal{S}}
\newcommand{\cY}{\mathcal{Y}}
\newcommand{\cn}[1]{N_\de({#1})}
\newcommand{\de}{\delta}
\newcommand{\e}{\epsilon}
\title{A note on the sum-product problem for fractal sets}
\author{Adam Cushman, William O'Regan}
\date{Source: arXiv:2604.21949v1 (CC BY 4.0)}
\begin{document}
\maketitle

\noindent\emph{Source and licence.} A note on the sum-product problem for fractal sets. Version arXiv:2604.21949v1 (CC BY 4.0), distributed by arXiv under the Creative Commons Attribution 4.0 International licence, http://creativecommons.org/licenses/by/4.0/, as recorded in the arXiv licence metadata for this exact version and shown on the abstract page https://arxiv.org/abs/2604.21949v1. Full licence notice: \texttt{licenses/arxiv-2604.21949-CC-BY-4.0.txt}.\par\smallskip

\noindent\textit{Editorial scope.} Counted unit: the article's lemma lem.sumlowerbounds (source0.tex lines 864--869). The article's complete original proof of it is delivered with it (source0.tex lines 870--1101, one \texttt{proof} environment of 232 lines). No mathematics is written, repaired or re-derived: the statement and the proof are the article's own lines, in the article's own notation, and no retained line is substituted or re-flowed.  Retained uncounted as the source-local context the proof needs: lines 387--389; lines 506--511; lines 784--793.  Nothing was omitted from any retained span, so the package carries no omission record.  No retained line contains a citation, so the wrapper carries no bibliography and no bibliography entry.  No retained line contains a figure environment, an image-inclusion command or drawing code, so no image is delivered and the package declares no asset dependency.  Editorial formatting only, in addition: an AMS \texttt{amsart} preamble that restates the article's own declarations and macros used by the retained lines, namely the environments \texttt{lemma}, \texttt{proposition} on separate counters, together with the standard packages those lines need.  The retained environments print in this document as Lemma 1, Proposition 1 in order of appearance, whereas the article prints the same items with the numbering of its own full document.  The complete native source of the article as distributed in the arXiv e-print for this exact version is bundled unchanged as \texttt{sources/199061/source0.tex}.
	\begin{lemma}\label{lem.unifsubset}
		Let $\epsilon > 0$ and suppose that $\delta >0 $ is small enough in terms of $\e.$ Let $A \subset \R^d$ be a $\delta$-separated set. Then there is an $\e$-uniform $A_0 \subset A$ with $\#A_0 \gtrsim \delta^{\epsilon}\#A.$ 
	\end{lemma}

\begin{lemma}\label{lem.energybounds}
  Let \(A \subset [1 / 2 , 1]\) be \((\delta,s,C_1)\)-KT and \(B \subset [-1,1]\) be \((\delta,t,C_2)\)-KT. Suppose that \(2s \leq 2 - t\). We have
\[
    \mathrm{E} _{3,\delta} (A,B) \lesssim  \delta^{- O(\epsilon)} C_1^{3}  C_2 ^{2}\frac{ \delta^{-s - t} N_{\delta} (AA)^{2} \# B}{\# A ^{2}} 
.\]
\end{lemma}

\begin{proposition}\label{prop:e2-cauchy-schwarz}
For any finite \(A,B \subset \mathbb{R} \),
\[
	\mathrm{E} _{2,\delta} (A,B) \sim  \mathrm{E} _{2,\delta} (A,-B)
,\]
and as a consequence,
\[
	\mathrm{E} _{2,\delta} (A,B) \gtrsim \frac{\# A ^{2} \# B ^{2}}{N_{\delta} \left( A + B \right) } 
.\]
\end{proposition}

\begin{lemma}\label{lem.sumlowerbounds}
We have, for any \(\eta > 0\), 
\[
    \delta^{O(\epsilon)} \delta^{- \sigma / 5    } \# A ^{\frac{11}{2} - \eta }\lesssim_{\eta}  \delta^{- s /5} N_{\delta} (A+A) ^{\frac{11}{5} } N_{\delta} (AA)^{\frac{11}{5} }
.\]
\end{lemma}

\begin{proof}
By Proposition \ref{prop:e2-cauchy-schwarz}, we have
\begin{equation} \label{eq:rr-lower-bound}
	\mathrm{E} _{2,\delta} (R_1,R_2) \gtrsim \frac{\# R_1 ^{2} \# R_2 ^{2}}{\cn{R_1 + R_2}} \gtrsim \frac{\delta^{4\epsilon} \# A ^{4}}{\# S},
\end{equation}
where we have used \(R_{1} + R_{2} \subset A + A\) and \(\# S = N_{\delta} (A+A)\). Thus, it remains to find
an upper bound on \(\mathrm{E} _{2,\delta} (R_1,R_2)\). It will be convenient to first find an
upper bound on \(\mathrm{E} _{\frac{7}{4}  ,\delta} (R_1,R_2)\), and then interpolate this with \(\mathrm{E} _{3,\delta} (A)\).

We use dyadic pigeonholing to find a level set of differences \(R_1-R_2\) supporting the \(7 / 4 \)-energy.
Namely, we write
\begin{align*}
	\mathrm{E} _{\frac{7}{4}  ,\delta} (R_1,R_2) & \sim \sum _{I \in \mathcal{D} _{\delta} (R_1-R_2)} \# \left( \pi ^{-1} (I) \cap (R_1 \times R_2)  \right) ^{\frac{7}{4}  }\\
	& = \sum _{\Delta \text{ dyadic} } \sum _{I \in \mathcal{P} _{\Delta} } \# \left( \pi ^{-1} (I) \cap (R_1 \times R_2) \right) ^{\frac{7}{4}  },
\end{align*}
where
\[
	\mathcal{P} _{\Delta} = \left\{ I \in \mathcal{D} _{\delta} (R_1 - R_2) : \# \left( \pi ^{-1} (I) \cap \left( R_1 \times R_2 \right)  \right) 	\in [\Delta,2\Delta) \right\} 
.\]	
Noting that for any interval \(I\),
\[
	\# \left( \pi ^{-1} (I) \cap (R_1 \times R_2) \right) \leq \# A ^{2}
,\]
it follows that there are
\[
	\lesssim \log \# A \lesssim _{\eta} \# A ^{\eta}
\]
many dyadic levels, and hence there is \(\Delta \in \mathbb{R} \) for which
\[
	\mathrm{E} _{\frac{7}{4} ,\delta} (R_1,R_2) \lesssim _{\eta} \# A ^{\eta} \sum _{I \in \mathcal{P} _{\Delta} } \#\left( \pi ^{-1} (I) \cap (R_1 \times R_2) \right) ^{\frac{7}{4}  } \sim \# A^{\eta} \Delta ^{\frac{7}{4}  } \# \mathcal{P} _{\Delta} 
.\]


Call \(P_{\Delta} = D \cap \bigcup _{I \in \mathcal{P} _{\Delta} } I\), where we recall that \(D\) is \(\delta\)-separated with \(\# D = N_{\delta} (A-A)\). 
We refine further appealing to Lemma \ref{lem.unifsubset} for the existence of
an \(\epsilon\)-uniform \(P_{\Delta} '\), with \(\# P_{\Delta} '\gtrsim \delta^{O(\epsilon)} \# P_{\Delta}  \).
Using the \(\epsilon\)-uniformity, we find that \(P_{\Delta} '\) is fit for energy estimates,
as

\begin{proposition}\label{prop:pdeltaKT}
\(P_{\Delta} '\) is a \((\delta ,s , \delta^{-O (\epsilon)} \# A / \Delta)\)-KT set. 
\end{proposition}
\begin{proof}
Let \(\mathcal{P} _{\Delta} '\) be those intervals of \(\mathcal{P} _{\Delta} \) which have intersection
with \(P_{\Delta} '\). We have
\[
	\sum _{I \in \mathcal{P} _{\Delta} '}  \#  \left( \pi ^{-1} (I) \cap A \times A \right) \geq  \sum _{I \in \mathcal{P} _{\Delta} '} \#  \left( \pi ^{-1} (I) \cap R_1 \times R_2 \right) \thicksim \Delta \# \mathcal{P} _{\Delta} '
,\]
so letting \(R_0\) be the rich set
\[
	R_0 = \left\{ x \in A  : \# \left( (x - A) \cap \mathcal{N} _{\delta} (P_{\Delta} ') \right) \gtrsim \frac{\Delta \# \mathcal{P} _{\Delta} '}{\# A} \right\}  
,\]
we have
\[
	\# R_0 \gtrsim \frac{\Delta \# \mathcal{P} _{\Delta}' }{\# A}
.\]
Since \(R_0\) is nonempty, let \(x_0 \in R_0\) be such that
\[
	\# \left( (x_0-A)\cap N_{\delta} (P_{\Delta} ') \right) \gtrsim \frac{\Delta \# \mathcal{P} _{\Delta} '}{\# A } \gtrsim \delta^{O(\epsilon)}\frac{\Delta \# P_{\Delta} }{\# A}    
.\]

Since \(P_{\Delta} '\) is \(\epsilon\)-uniform and \(\delta\)-separated, for any \(\delta \leq \rho \leq 1\),
\[
	N_{\delta} \left( P_{\Delta} ' \cap B(x,\rho) \right) N_{\rho} (P_{\Delta} ') \thicksim \# P_{\Delta} '
.\]
Note that
\[
	N_{\rho} (P_{\Delta} ') \gtrsim  N_{\rho} \left( (x_0 - A)\cap \mathcal{N} _{\delta} (P_{\Delta} ') \right) 
,\]
and since \(x_0 - A\) is a \(\delta\)-separated \((\delta,s, \delta^{-O(\epsilon)})\)-KT set, the number of
points of \((x_0 - A)\cap \mathcal{N} _{\delta} (P_{\Delta} ') \) in a \(\rho\) ball is
\[
	(x_0 - A)\cap \mathcal{N} _{\delta} (P_{\Delta} ')\cap B(x,\rho)  \lesssim \delta^{-O(\epsilon)} \left( \frac{\rho}{\delta}  \right) ^{s}
,\]
and hence the number of such balls is 
\[
	N_{\rho} \left( (x_0 - A)\cap \mathcal{N} _{\delta} (P_{\Delta} ') \right) \gtrsim \frac{\# \left( x_0 - A \right) \cap \mathcal{N} _{\delta} (P_{\Delta} ')}{\delta^{-O(\epsilon)} \rho^{s} \delta^{-s}}    \gtrsim \delta^{O(\epsilon)} \left( \delta^{s} \frac{\Delta \# P_{\Delta} }{\# A}  \right) \rho ^{-s}
,\]
and hence
\[
	N_{\delta} \left( P_{\Delta} ' \cap B(x,\rho) \right) \lesssim \frac{\# P_{\Delta} '}{N_{\rho} (P_{\Delta} ')   } \lesssim \delta^{- O(\epsilon)} \frac{\# A}{\Delta}  \left( \frac{\rho}{\delta}  \right) ^{s}
.\] 
\end{proof}
Because \(- P_{\Delta} '\) is a \((\delta,s,\delta^{-O(\epsilon)} \# A / \Delta)\)-KT set, we may apply Lemma \ref{lem.energybounds} to give
\[
	\mathrm{E} _{3,\delta} (A,- P_{\Delta} ')\lesssim \delta^{-O(\epsilon)} \frac{\delta^{-2s} N_{\delta} (AA)^{2} \# P_{\Delta} }{\Delta^{2}} 
.\]
We will use this momentarily.
We proceed now as we did in the case of the difference-set, the change being that now we use the
truism
\[
	(r_1 + a) - (r_2 + a) = r_1 - r_2
,\]
to find solutions to the equation \(p_1 - p_2 = p_3 + O(\delta)\) where \(p_1 \in \mathcal{N} _{\delta} (S_{\text{pop}} ) \), \(p_2 \in \mathcal{N} _{\delta} (P_{+}) \), \(p_3 \in \mathcal{N} _{\delta} (P_{\Delta}' )\). 

Carrying this out, we map the set of triples
\[
	X = \left\{ (r_1,r_2,a)\in R_1 \times R_2 \times A : r_1 + a \in \mathcal{N} _{\delta} (S_{\text{pop}} ),~ r_2 + a \in \mathcal{N} _{\delta} (P_{+} ),~ r_1 - r_2 \in \mathcal{N} _{\delta} (P_{\Delta}' ) \right\} 
\]
by
\[
	f : (r_1,r_2,a) \mapsto (r_1 + a ,~ r_2 + a)
.\]
Let \(\mathcal{Y} \) be those intervals into which \(f\) maps \(X\), that is
\[
	\mathcal{Y} = \left\{ (I,J)\in\cS \times \cP_{+}  	: \left( I - J \right) \cap \mathcal{N} _{\delta} (P_{\Delta} ') \neq \varnothing  \right\} 	
,\]
where we note that
\[
	f(X) \subset \bigcup_{B \in \cY} B
.\]

We first obtain a lower bound on \(\# X\).
We choose the rich elements \((r_1,r_2)\) to be all those whose difference is in our dyadic level.
By the definition of \(\mathcal{P} _{\Delta} '\), 
\[
	\Delta \# \mathcal{P} _{\Delta}' \leq \sum _{I \in \mathcal{P} _{\Delta}' } \# \left( \pi ^{-1} (I)\cap (R_1 \times R_2) \right) \leq  \# \left\{ (r_1,r_2)\in R_1 \times R_2 : r_1 - r_2 \in \mathcal{N} _{\delta} (P_{\Delta}' )  \right\} 
.\]
We then choose all the \(a\) to be those which make \(r_1 + a \in \mathcal{N} _{\delta} (S_{\text{pop}} ) \) and \(r_2 + a \in \mathcal{N} _{\delta} (P_{+} 	) \).
For a fixed pair \((r_1,r_2)\in R_1 \times R_2\) by the definition of \(R_1\),
\[
	\# \left\{ a \in A : r_1 + a \in \mathcal{N} _{\delta} (S_{\text{pop}} ) \right\} \geq c \cdot \delta^{\epsilon}\# A
.\]
Call the set of all such \(a\) to be \(A_1\). Then, by the definition of \(R_2\),
\[
	\# \left\{ a \in A : r_2 + a \in \mathcal{N} _{\delta} (P_{+} )  \right\} \geq \left( 1 - \frac{c \delta ^{\epsilon}}{2}  \right) \# A
.\]
Call the set of all such \(a\) to be \(A_2\).
Using inclusion-exclusion,
\[
	\# \left( A_1 \cap A_2 \right) = \# A_1 + \# A_2 - \# \left( A_1 \cup A_2 \right) 
,\]
and hence
\[
	\# \left( A_1 \cap A_2 \right) \geq c \delta^{\epsilon}\# A + \# A \left( 1 - \frac{c \delta^{\epsilon}}{2}  \right) - \# A \gtrsim \delta^{\epsilon} \# A
.\]
Taking a union over all such \((r_1,r_2)\in R_1 \times R_2\), we find that
\[
	\# X \gtrsim \delta^{\epsilon} \Delta \# \mathcal{P} _{\Delta}' \# A \gtrsim \delta^{O(\epsilon)} \Delta \# \mathcal{P} _{\Delta}  \# A
.\]

By a nearly identical argument to that of the difference-set portion,
\[
	\# \left\{ (x_1,x_2)\in X^{2} : \left\lvert f(x_1) - f(x_2) \right\rvert \lesssim \delta \right\} \lesssim \mathrm{E} _{3,\delta} (A)
.\]
Thus, again by Cauchy-Schwarz, we obtain
\begin{equation} \label{eq:sum-post-CS}
	\delta^{O(\epsilon)} \Delta^{2} \# \mathcal{P} _{\Delta} ^{2}\# A ^{2} \lesssim \mathrm{E} _{3,\delta} (A) \# \mathcal{Y} .
\end{equation}

With \(\mathcal{Y} \) in terms of \(P_{\Delta} '\) instead of \(P_{\Delta} \), we may expand the \(\mathcal{P} _{+} \) term
and then use incidence estimates for the remaining terms.
We see that
\[
	\# \mathcal{Y}  \sim \sum _{p_1 \in S_{\text{pop}}   } \sum _{p_2 \in P_{+} } \sum _{p_3 \in P_{\Delta} ' }  1_{\left\lvert p_1 - p_2 - p_3 \right\rvert \leq \delta}
.\]
Expanding \(p_2\), we note that by the definition of \(\mathcal{P} _{+}  \) we have
\[
	\sum _{p_1 \in S_{\text{pop}} } \sum _{p_2 \in P_{+} } \sum _{p_3 \in P_{\Delta}' } \sum _{a_1,a_2 \in A} 1_{\left\lvert a_1 + a_2 - p_2 \right\rvert \leq \delta} 1_{\left\lvert p_1-p_2-p_3 \right\rvert \leq \delta} \gtrsim \frac{\# A ^{2} \delta^{\epsilon}}{\# S} \cdot \# \mathcal{Y} 
,\]
so eliminating the \(P_{+} \) term gives
\[
	\# \mathcal{Y}  \lesssim \frac{\# S}{\# A ^{2}\delta^{\epsilon}} \sum _{p_1 \in S_{\text{pop}} } \sum _{p_3 \in P_{\Delta}' } \sum _{a_1,a_2 \in A} 1_{\left\lvert p_1-p_3-a_1 - a_2 \right\rvert \leq \delta} 
.\]
We group this as \((a_1 - p_1) + (a_2 + p_3)\) to obtain that
\[
	\# \mathcal{Y}  \lesssim \frac{\# S}{\# A ^{2} \delta^{\epsilon}} 	\sum _{z \in \delta \mathbb{Z} } \left( \sum _{a_2 \in A} \sum _{p_3 \in P_{\Delta}' }  1_{\left\lvert a_2 + p_3 + z \right\rvert \leq \delta}  \right) \left( \sum _{a_1 \in A} \sum _{p_1 \in S _{\text{pop}} 	}  1_{\left\lvert a_1 - p_1 - z \right\rvert \leq \delta} \right) 
,\]
and hence by using H\"older we get
\[
	\# \mathcal{Y}  \lesssim \frac{\# S}{\# A ^{2}\delta^{\epsilon}} \mathrm{E} _{\frac{3}{2} ,\delta} (A,- P_{\Delta}' )^{\frac{2}{3} } \mathrm{E} _{3,\delta} (A, S_{\text{pop}} )^{\frac{1}{3} }
.\]

Combining, we have
\[
	\frac{\Delta ^{2} \# P_{\Delta} ^{2} \# A ^{4} }{\# S} \lesssim \delta^{- O(\epsilon)}\mathrm{E} _{3,\delta} (A) \mathrm{E} _{3,\delta} (A,S_{\text{pop}} )^{\frac{1}{3} } \mathrm{E} _{\frac{3}{2}  ,\delta} (A,P_{\Delta} ')^{\frac{2}{3} }
.\]
Recall our energy bound
\[
	\mathrm{E} _{3,\delta} (A,- P_{\Delta} ')\lesssim \delta^{-O(\epsilon)} \frac{\delta^{-2s} N_{\delta} (AA)^{2} \# P_{\Delta} }{\Delta^{2}} 
.\]
Interpolating for the \(3/2\) energy using this, we see by H\"older that
\[
	\mathrm{E} _{\frac{3}{2} ,\delta} (A, - P_{\Delta} ') \leq \mathrm{E} _{1,\delta} (A , - P_{\Delta} ')^{\frac{3}{4} } \mathrm{E} _{3,\delta} (A , - P_{\Delta} ')^{\frac{1}{4} }
,\]
which upon substituting gives
\[
	\mathrm{E} _{\frac{3}{2} ,\delta} (A, -P_{\Delta} ')^{\frac{2}{3} } \lesssim \frac{ \delta^{- O(\epsilon)} \# A^{\frac{1}{2} } \# P_{\Delta} ^{\frac{2}{3} } \delta^{- s / 3} N_{\delta} (AA)^{\frac{1}{3} }}{\Delta^{\frac{1}{3} }} 
.\]
Using the bound on \(\mathrm{E} _{\frac{3}{2} , \delta} (A, - P_{\Delta} ')\) along with
\[
	\mathrm{E}_{3, \delta}\left(A, S_{\mathrm{pop}}\right) \lesssim \delta^{-O(\epsilon)} \frac{N_\delta(A A)^2 N_\delta\left(S_{\mathrm{pop}}\right)^3}{\# A \delta^{-\sigma}}
\]
and
\[
	\mathrm{E}_{3, \delta}(A) \lesssim \delta^{-O(\epsilon)} N_\delta(A A)^2 \# A,
\]
gives
\[
	\frac{\Delta ^{2} \# P_{\Delta} ^{2} \# A ^{4}}{\# S} \lesssim \delta^{- O(\epsilon)} \left( N_\delta(A A)^2 \# A \right) \left( \frac{N_\delta(A A)^2 N_\delta\left(S_{\mathrm{pop}}\right)^3}{\# A \delta^{-\sigma}} \right) ^{\frac{1}{3} }  \frac{ \# A^{\frac{1}{2} } \# P_{\Delta} ^{\frac{2}{3} } \delta^{- s / 3} N_{\delta} (AA)^{\frac{1}{3} }}{\Delta^{\frac{1}{3} }} 
\]
which upon simplifying is
\[
	\Delta ^{\frac{7}{3} } \# P_{\Delta} ^{\frac{4}{3} }  \lesssim \# S\# A ^{- \frac{17}{6} }  N_{\delta} (AA)^{3} N_{\delta} (S_{\text{pop}} )   \delta^{ - s / 3 + \sigma / 3   }
.\]

We have chosen our level \(\Delta\) so that
\[
	\Delta^{\frac{7}{3} } \# P_{\Delta} ^{\frac{4}{3} }= \left( \Delta^{\frac{7}{4} } \# P_{\Delta} \right) ^{\frac{4}{3} } \gtrsim _{\eta} \# A ^{- \eta} \mathrm{E} _{\frac{7}{4} ,\delta} (R_1,R_2)^{\frac{4}{3} }   
.\]
Interpoling for the second energy, we have by H\"older that
\[
	\mathrm{E} _{2,\delta}  (R_1,R_2) \leq \mathrm{E} _{\frac{7}{4} ,\delta} (R_1,R_2)^{\frac{4}{5} } \mathrm{E} _{3,\delta} (A)^{\frac{1}{5} }
,\]
so substituting gives
\[
	\mathrm{E} _{2,\delta}  (R_1,R_2) \lesssim_{\eta}  \# A ^{\eta} \left[  \# S\# A ^{- \frac{17}{6} }  N_{\delta} (AA)^{3} N_{\delta} (S_{\text{pop}} )   \delta^{ - s / 3 + \sigma / 3   } \right] ^{\frac{3}{5}  } \left( \delta^{-O(\epsilon)} N_\delta(A A)^2 \# A \right) ^{\frac{1}{5} }
\]
or by simplifying,
\[
	\mathrm{E} _{2,\delta}  (R_1,R_2) \lesssim_{\eta} \# A^{\eta} \delta^{- O(\epsilon)} \frac{\# S ^{\frac{3}{5} } N_{\delta} (AA) ^{\frac{11}{5} } N_{\delta} (S_{\text{pop}} )^{\frac{3}{5} } \delta ^{- s /5 + \sigma /5}}{\# A ^{\frac{3}{2} }} 
.\]
Using our lower bound on the second energy, we have
\[
	\frac{\# A ^{4 - \eta}}{\# S} \lesssim_{\eta}  \delta^{-O(\epsilon)} \frac{\# S ^{\frac{3}{5} } N_{\delta} (AA) ^{\frac{11}{5} } N_{\delta} (S_{\text{pop}} )^{\frac{3}{5} } \delta ^{- s /5 + \sigma /5}}{\# A ^{\frac{3}{2} }} 
,\]
which simplifies to 
\[
	\delta^{O(\epsilon)} \delta^{- \sigma / 5    } \# A ^{\frac{11}{2} - \eta }\lesssim_{\eta}  \delta^{- s /5} N_{\delta} (A+A) ^{\frac{11}{5} } N_{\delta} (AA)^{\frac{11}{5} }
,\]
and we are finished.
\end{proof}

\end{document}
